\chapter{Fever}

\section*{Definition}
Fever is defined as an elevation of core body temperature above the normal diurnal range, typically exceeding 38.0\textdegree{}C (100.4\textdegree{}F) when measured orally or 38.3\textdegree{}C (100.9\textdegree{}F) rectally. It represents a regulated increase in the hypothalamic set-point mediated by endogenous pyrogens during immune activation, distinguishing it from hyperthermia where thermoregulation fails.

\section*{Pathophysiology}
Fever results from the action of exogenous and endogenous pyrogens on the hypothalamus. Exogenous pyrogens stimulate immune cells to release endogenous pyrogens including interleukin-1 (IL-1), IL-6, tumor necrosis factor-alpha (TNF-$alpha$), and interferon-alpha. These cytokines induce cyclooxygenase-2 expression leading to increased prostaglandin E2 synthesis which raises the thermoregulatory set-point. Heat-conservation mechanisms (vasoconstriction, shivering) activate until temperature reaches new set-point. Resolution involves active heat-loss (vasodilation, sweating).

\section*{Etiology}
\begin{itemize}
  \item Infectious causes (most common): Viral infections (URI, influenza, enteroviruses, hepatitis, HIV), Bacterial infections (pneumonia, UTI, cellulitis, endocarditis, abscesses, meningitis), Fungal infections (histoplasmosis, coccidioidomycosis), Parasitic infections (malaria, toxoplasmosis)
  \item Inflammatory/autoimmune conditions: Rheumatoid arthritis, systemic lupus erythematosus, vasculitides (giant cell arteritis, Takayasu arteritis), inflammatory bowel disease (Crohn disease, ulcerative colitis), sarcoidosis
  \item Neoplastic causes: Lymphomas (especially Hodgkin lymphoma), leukemias, renal cell carcinoma, hepatocellular carcinoma, melanoma, multiple myeloma
  \item Drug-induced fever: Antibiotics (sulfonaminals, beta-lactams), anticonvulsants (phenytoin, carbamazepine), antiarrhythmics (quinidine), biologics, vaccines
  \item Thromboembolic events: Pulmonary embolism, deep vein thrombophlebitis
  \item Other causes: Heat stroke, CNS hemorrhage, post-operative state, factitious disorder
\end{itemize}

\section*{Indications for Evaluation}
Seek care if:
\begin{itemize}
  \item Fever persisting beyond 3 days without identified source
  \item Temperature $\ge$ 40\textdegree{}C (104\textdegree{}F) regardless of duration
  \item Fever accompanied by neck stiffness, photophobia, altered mental status (meningeal irritation)
  \item Immunocompromised patient with any fever (temperature $\ge$ 38.3\textdegree{}C)
  \item Fever with rash that does not blanch (possible meningococcemia)
  \item Severe headache, abdominal pain, flank pain with fever
\end{itemize}

\section*{Related Symptoms}
\begin{itemize}
  \item Chills
  \item Rigors
  \item Diaphoresis
  \item Headache
  \item Myalgia
  \item Arthralgia
  \item Fatigue
  \item Anorexia
  \item Weight loss
  \item Night sweats
  \item Dehydration
  \item Tachycardia
  \item Tachypnea
\end{itemize}
\textbf{Note:} Persistent fever warrants systematic investigation including blood cultures, CBC with differential, inflammatory markers (ESR, CRP), urinalysis, chest imaging, and consideration of autoimmune/neoplastic etiologies when initial workup negative.

\section*{Self-Care / Home Management}
\begin{itemize}
  \item Rest and maintain adequate hydration with water, clear fluids, or oral rehydration solutions.
  \item Use acetaminophen or ibuprofen to reduce fever and body aches (follow label dosing instructions).
  \item Dress lightly and keep the room at a comfortable temperature — avoid bundling excessively.
  \item Monitor temperature at regular intervals and track other symptoms to identify worsening trends.
  \item Seek medical attention for fever above 40	extdegree{}C (104	extdegree{}F), fever lasting more than 3 days, or fever with stiff neck, severe headache, rash, or difficulty breathing.
\end{itemize}

\section*{Diagnostic Approach}
\begin{itemize}
  \item History includes fever duration, pattern (intermittent, remittent, continuous), associated symptoms (chills, rigors, sweats), and exposure history.
  \item Vital signs assess fever severity and hemodynamic stability; complete physical examination seeks a source.
  \item Laboratory evaluation includes CBC with differential, blood cultures, urinalysis, inflammatory markers (CRP, ESR), and targeted testing based on exposure history.
  \item Imaging (chest X-ray, abdominal ultrasound, CT) is guided by clinical findings and suspected source.
\end{itemize}

\section*{Treatment / Management Overview}
\begin{itemize}
  \item Antipyretics (acetaminophen, NSAIDs) provide symptomatic relief but do not treat the underlying cause.
  \item Empiric antibiotics may be indicated for suspected bacterial infection after appropriate cultures are obtained.
  \item Antiviral, antifungal, or antimalarial therapy is directed at specific pathogens based on exposure and testing.
  \item Source control (drainage of abscess, removal of infected device) is essential for focal infections.
\end{itemize}

\clearpage
